% xelatex graph_analyzer_logic.tex
\documentclass[10pt]{article}
\usepackage{fontspec}
\usepackage{polyglossia}
\setdefaultlanguage{russian}
\setmainfont{Liberation Serif}
\setsansfont{Liberation Sans}
\setmonofont[Scale=0.80]{DejaVu Sans Mono}

\usepackage[
  paperwidth=297mm, paperheight=210mm,  % A4 landscape
  top=14mm, bottom=14mm,
  left=14mm, right=14mm,
]{geometry}

\usepackage[dvipsnames,svgnames,table]{xcolor}
\usepackage{tikz}
\usetikzlibrary{
  arrows.meta, positioning, fit,
  shapes.geometric, shapes.multipart,
  backgrounds, calc, shadows.blur,
}
\usepackage[unicode,hidelinks]{hyperref}

% ── палитра ──────────────────────────────────────────────────────
\definecolor{CInput}  {RGB}{30,  80, 160}   % входные данные
\definecolor{CParse}  {RGB}{90,  40, 160}   % парсинг
\definecolor{CAlgo}   {RGB}{140, 20, 100}   % алгоритмы
\definecolor{CPrompt} {RGB}{15, 130,  90}   % промпты
\definecolor{CLLM}    {RGB}{20, 160, 110}   % LLM
\definecolor{CMerge}  {RGB}{180, 90,  10}   % merge
\definecolor{COut}    {RGB}{20,  80,  20}   % вывод

\definecolor{BInput}  {RGB}{225,235,255}
\definecolor{BParse}  {RGB}{240,230,255}
\definecolor{BAlgo}   {RGB}{255,225,245}
\definecolor{BPrompt} {RGB}{210,250,235}
\definecolor{BLLM}    {RGB}{195,248,225}
\definecolor{BMerge}  {RGB}{255,238,210}
\definecolor{BOut}    {RGB}{220,245,220}

% ── стили ────────────────────────────────────────────────────────
\tikzset{
  % процессный блок
  proc/.style={
    rectangle, rounded corners=4pt,
    minimum width=38mm, minimum height=11mm,
    text centered, draw, line width=0.7pt,
    font=\sffamily\footnotesize,
    blur shadow={shadow blur steps=3, shadow xshift=0.4mm,
                 shadow yshift=-0.5mm, shadow blur radius=1.2mm,
                 shadow opacity=20},
  },
  % блок с данными — простой прямоугольник с разделителем
  data/.style={
    rectangle, rounded corners=3pt,
    minimum width=32mm, text width=30mm,
    draw=#1!55, fill=#1!12, line width=0.6pt,
    font=\sffamily\scriptsize, align=left,
    inner sep=4pt,
    blur shadow={shadow blur steps=3, shadow xshift=0.3mm,
                 shadow yshift=-0.4mm, shadow blur radius=1mm,
                 shadow opacity=12},
  },
  % стрелки
  arr/.style={-{Stealth[length=4.5pt,width=3.5pt]}, line width=0.9pt},
  arr-fat/.style={arr, line width=1.3pt},
  lbl/.style={font=\sffamily\tiny, fill=white, inner sep=1.5pt, midway},
  % аннотации типов
  typ/.style={font=\ttfamily\tiny, color=gray!60, inner sep=0pt},
}

\pagestyle{empty}

% ════════════════════════════════════════════════════════════════
\begin{document}
\begin{center}
{\large\bfseries\sffamily Граф-аналитик процессов — логика потока}\\[-1pt]
{\scriptsize\color{gray} поток-граф-аналитик · Common Lisp · data-flow}
\end{center}

\vspace{-2mm}

\begin{center}
\scalebox{0.58}{%
\begin{tikzpicture}[
  node distance=6mm and 10mm,
  x=1mm, y=1mm,
]

%% ══════════════════════════════════════════════
%% КОЛОНКА 0 — Входные данные
%% ══════════════════════════════════════════════

\node[proc, fill=BInput, draw=CInput!60, text=CInput!90!black,
      minimum width=36mm, minimum height=22mm, text width=34mm,
      align=center] (in-graph)
  {\textbf{Mermaid flowchart}\\[2pt]
   {\ttfamily\tiny граф.md}\\
   {\tiny текстовый граф:\\flowchart TD\\A0[…] –{}–> B1[…]}};

\node[proc, fill=BInput, draw=CInput!60, text=CInput!90!black,
      minimum width=36mm, minimum height=22mm, text width=34mm,
      align=center, below=8mm of in-graph] (in-models)
  {\textbf{Django-модели}\\[2pt]
   {\ttfamily\tiny модели.py}\\
   {\tiny class WasteSite…:\\  status = models.FK…}};

%% ══════════════════════════════════════════════
%% КОЛОНКА 1 — Парсинг
%% ══════════════════════════════════════════════

\node[proc, fill=BParse, draw=CParse!60, text=CParse!90!black,
      right=16mm of in-graph, minimum width=36mm, text width=34mm,
      align=center] (parse-graph)
  {\textbf{\ttfamily разобрать-граф}\\[2pt]
   {\tiny regex (cl-ppcre):\\
   {\ttfamily (\textbackslash w+)["](…)["]}\\
   {\ttfamily (\textbackslash w+)-->|"…"|}\\
   → classify узлы}};

\node[proc, fill=BParse, draw=CParse!60, text=CParse!90!black,
      right=16mm of in-models, minimum width=36mm, text width=34mm,
      align=center] (parse-models)
  {\textbf{\ttfamily разобрать-модели}\\[2pt]
   {\tiny regex по строкам:\\
   {\ttfamily class (\textbackslash w+)\textbackslash s*(}\\
   {\ttfamily (\textbackslash w+)=models.(\textbackslash w+)(}\\
   → choices? null?}};

%% структуры данных после парсинга
\node[data=CParse, right=12mm of parse-graph] (ds-graph)
  {{\color{CParse!80!black}\textbf{узел}}\\[2pt]
   {\ttfamily\tiny id: string\\метка: string\\ошибка?: bool\\терминал?: bool\\решение?: bool}};

\node[data=CParse, right=12mm of parse-models] (ds-models)
  {{\color{CParse!80!black}\textbf{модель}}\\[2pt]
   {\ttfamily\tiny имя: string\\поля: [поле]\\\ \ .имя .тип\\\ \ .null? .choices?}};

% ── вспомогательно: ребро ────────────────────────────────────────
\node[data=CParse, below=5mm of ds-graph] (ds-edge)
  {{\color{CParse!80!black}\textbf{ребро}}\\[2pt]
   {\ttfamily\tiny откуда: string\\куда: string\\условие: string?}};

%% ══════════════════════════════════════════════
%% КОЛОНКА 2 — Алгоритмы подготовки
%% ══════════════════════════════════════════════

\node[proc, fill=BAlgo, draw=CAlgo!50, text=CAlgo!90!black,
      right=28mm of ds-graph, yshift=-2mm,
      minimum width=40mm, text width=38mm, align=center] (algo-bfs-fwd)
  {\textbf{\ttfamily недостижимые}\\[2pt]
   {\tiny BFS вперёд от стартовых узлов\\
   (нет входящих рёбер)\\
   → множество непосещённых}};

\node[proc, fill=BAlgo, draw=CAlgo!50, text=CAlgo!90!black,
      below=7mm of algo-bfs-fwd,
      minimum width=40mm, text width=38mm, align=center] (algo-bfs-bwd)
  {\textbf{\ttfamily пути-к-ошибкам}\\[2pt]
   {\tiny BFS назад от ERR-узлов\\
   глубина 4 → цепочка\\
   предшествующих шагов}};

\node[proc, fill=BAlgo, draw=CAlgo!50, text=CAlgo!90!black,
      below=7mm of algo-bfs-bwd,
      minimum width=40mm, text width=38mm, align=center] (algo-summary)
  {\textbf{\ttfamily резюме-графа/моделей}\\[2pt]
   {\tiny текстовое описание\\
   для LLM-контекста:\\
   «узлов N, рёбер M…»}};

%% выходы алгоритмов
\node[data=CAlgo, right=10mm of algo-bfs-fwd, yshift=3mm] (ds-unreach)
  {{\color{CAlgo!80!black}\textbf{недост.}}\\[2pt]
   {\ttfamily\tiny [узел]\\список изолированных}};

\node[data=CAlgo, right=10mm of algo-bfs-bwd] (ds-errpaths)
  {{\color{CAlgo!80!black}\textbf{ошибки}}\\[2pt]
   {\ttfamily\tiny ERR1 ← E20, E19…\\строка для промпта}};

\node[data=CAlgo, right=10mm of algo-summary, yshift=-3mm] (ds-summary)
  {{\color{CAlgo!80!black}\textbf{резюме}}\\[2pt]
   {\ttfamily\tiny граф-строка\\модели-строка}};

%% ══════════════════════════════════════════════
%% КОЛОНКА 3 — Три LLM-промпта
%% ══════════════════════════════════════════════

\node[proc, fill=BPrompt, draw=CPrompt!60, text=CPrompt!90!black,
      right=22mm of ds-unreach, yshift=8mm,
      minimum width=42mm, text width=40mm, align=center] (p1)
  {\textbf{① :компенсация}\\[2pt]
   {\tiny граф + модели + \underline{ошибки}\\
   «для каждого ERR-узла:\\что уже изменилось в БД\\и не откатывается?»}};

\node[proc, fill=BPrompt, draw=CPrompt!60, text=CPrompt!90!black,
      below=7mm of p1,
      minimum width=42mm, text width=40mm, align=center] (p2)
  {\textbf{② :крайние-случаи}\\[2pt]
   {\tiny граф + \underline{модели}\\
   «пустые списки, null FK,\\параллельные вызовы,\\partial execution»}};

\node[proc, fill=BPrompt, draw=CPrompt!60, text=CPrompt!90!black,
      below=7mm of p2,
      minimum width=42mm, text width=40mm, align=center] (p3)
  {\textbf{③ :достижимость}\\[2pt]
   {\tiny граф + \underline{недостижимые}\\
   «условия невозможные\\при нормальном потоке,\\незавершённые циклы»}};

%% LLM-вызов
\node[proc, fill=BLLM, draw=CLLM!60, text=CLLM!80!black,
      right=12mm of p2,
      minimum width=36mm, text width=34mm, align=center,
      minimum height=48mm] (llm)
  {\textbf{вызвать-llm}\\[3pt]
   {\ttfamily\tiny dexador:post\\OpenRouter API}\\[4pt]
   {\tiny \texttt{POST /chat/completions}\\
   model: claude-opus-4-5\\
   max\_tokens: 2000}\\[4pt]
   {\scriptsize ↓}\\[2pt]
   {\tiny JSON-ответ:\\
   \texttt{\{findings:[…]\}}}};

%% ══════════════════════════════════════════════
%% КОЛОНКА 4 — Merge + Format
%% ══════════════════════════════════════════════

\node[proc, fill=BMerge, draw=CMerge!60, text=CMerge!90!black,
      right=12mm of llm, yshift=15mm,
      minimum width=38mm, text width=36mm, align=center] (merge)
  {\textbf{\ttfamily извлечь-findings}\\[2pt]
   {\tiny parse JSON-ответа (cl-json)\\strip markdown-обёртки\\→ список alist}};

\node[data=CMerge, below=6mm of merge] (ds-finding)
  {{\color{CMerge!80!black}\textbf{finding}}\\[2pt]
   {\ttfamily\tiny node: "ERR1"\\category: "comp\_gap"\\description: "…"\\recommendation: "…"\\severity: critical}};

\node[proc, fill=BMerge, draw=CMerge!60, text=CMerge!90!black,
      below=6mm of ds-finding,
      minimum width=38mm, text width=36mm, align=center] (fmt)
  {\textbf{\ttfamily отформатировать-отчёт}\\[2pt]
   {\tiny сгруппировать по severity\\добавить иконки 🔴 ⚠️  ℹ️\\header со счётчиками}};

%% ══════════════════════════════════════════════
%% КОЛОНКА 5 — Выход
%% ══════════════════════════════════════════════

\node[proc, fill=BOut, draw=COut!60, text=COut!80!black,
      right=12mm of merge, yshift=-22mm,
      minimum width=36mm, text width=34mm, align=center,
      minimum height=40mm] (out)
  {\textbf{Отчёт}\\[4pt]
   {\scriptsize
   🔴 [Компенсация] ERR1\\
   \hspace{2mm}{\tiny не откатывает B2-B3}\\[3pt]
   ⚠️\ [Крайний случай] E5\\
   \hspace{2mm}{\tiny containers=[ ] }\\[3pt]
   🟡 [Достижимость] F2\\
   \hspace{2mm}{\tiny недостижим из E25}\\[4pt]
   {\tiny Всего: N  🔴 K  ⚠️\ M}}};

%% ══════════════════════════════════════════════
%% СТРЕЛКИ
%% ══════════════════════════════════════════════

% входы → парсеры
\draw[arr, color=CInput] (in-graph)  -- node[lbl]{текст} (parse-graph);
\draw[arr, color=CInput] (in-models) -- node[lbl]{текст} (parse-models);

% парсеры → структуры данных
\draw[arr, color=CParse] (parse-graph)  -- (ds-graph);
\draw[arr, color=CParse] (parse-graph)  -- (ds-edge);
\draw[arr, color=CParse] (parse-models) -- (ds-models);

% структуры → алгоритмы
\draw[arr, color=CParse!70] (ds-graph) -- node[lbl,above]{узлы+рёбра} (algo-bfs-fwd);
\draw[arr, color=CParse!70] (ds-graph) |- (algo-bfs-bwd);
\draw[arr, color=CParse!70] (ds-edge)  |- (algo-bfs-bwd);
\draw[arr, color=CParse!70] (ds-graph) |- (algo-summary);
\draw[arr, color=CParse!70] (ds-models) -| (algo-summary);

% алгоритмы → выходы
\draw[arr, color=CAlgo] (algo-bfs-fwd) -- (ds-unreach);
\draw[arr, color=CAlgo] (algo-bfs-bwd) -- (ds-errpaths);
\draw[arr, color=CAlgo] (algo-summary) -- (ds-summary);

% выходы → промпты (заполнение контекста)
\draw[arr, color=CAlgo!60] (ds-errpaths) -| node[lbl,near start]{} (p1);
\draw[arr, color=CAlgo!60] (ds-summary) -| (p1);
\draw[arr, color=CAlgo!60] (ds-summary) -- (p2);
\draw[arr, color=CAlgo!60] (ds-unreach) -| node[lbl,near start]{} (p3);
\draw[arr, color=CAlgo!60] (ds-summary) -| (p3);

% промпты → LLM (3 стрелки входят)
\draw[arr, color=CPrompt] (p1.east) -- (p1.east -| llm.west);
\draw[arr, color=CPrompt] (p2.east) -- (p2.east -| llm.west);
\draw[arr, color=CPrompt] (p3.east) -- (p3.east -| llm.west);

% LLM → merge (3 стрелки выходят)
\draw[arr-fat, color=CLLM] (llm.east |- p1.east) -- node[lbl,above]{\tiny findings} (merge.west |- p1.east);
\draw[arr-fat, color=CLLM] (llm.east |- p2.east) -- (merge.west |- p2.east);
\draw[arr-fat, color=CLLM] (llm.east |- p3.east) -- (merge.west |- p3.east);

% merge → ds-finding → fmt → out
\draw[arr, color=CMerge] (merge) -- (ds-finding);
\draw[arr, color=CMerge] (ds-finding) -- (fmt);
\draw[arr-fat, color=COut] (fmt.east) -| (out.south);

%% ══════════════════════════════════════════════
%% ФОНОВЫЕ ЗОНЫ (lanes)
%% ══════════════════════════════════════════════

\begin{scope}[on background layer]
  \node[fit=(in-graph)(in-models),
        fill=CInput!4, draw=CInput!20, rounded corners=5pt,
        inner sep=4pt, label={[font=\sffamily\tiny\color{CInput}]above:Входные данные}] {};

  \node[fit=(parse-graph)(parse-models)(ds-graph)(ds-edge)(ds-models),
        fill=CParse!4, draw=CParse!20, rounded corners=5pt,
        inner sep=5pt, label={[font=\sffamily\tiny\color{CParse}]above:Парсинг}] {};

  \node[fit=(algo-bfs-fwd)(algo-bfs-bwd)(algo-summary)(ds-unreach)(ds-errpaths)(ds-summary),
        fill=CAlgo!4, draw=CAlgo!20, rounded corners=5pt,
        inner sep=5pt, label={[font=\sffamily\tiny\color{CAlgo}]above:Подготовка контекста}] {};

  \node[fit=(p1)(p2)(p3)(llm),
        fill=CPrompt!4, draw=CPrompt!20, rounded corners=5pt,
        inner sep=5pt, label={[font=\sffamily\tiny\color{CPrompt}]above:LLM-анализ}] {};

  \node[fit=(merge)(ds-finding)(fmt),
        fill=CMerge!4, draw=CMerge!20, rounded corners=5pt,
        inner sep=5pt, label={[font=\sffamily\tiny\color{CMerge}]above:Объединение}] {};

  \node[fit=(out),
        fill=COut!4, draw=COut!20, rounded corners=5pt,
        inner sep=5pt, label={[font=\sffamily\tiny\color{COut}]above:Результат}] {};
\end{scope}

\end{tikzpicture}%
}% scalebox
\end{center}

\vfill
{\centering\scriptsize\color{gray}
  \texttt{поток-граф-аналитик} · cl-ppcre · dexador · cl-json ·
  3 специализированных промпта · категории: \texttt{compensation\_gap | edge\_case | unreachable\_branch}\par}

\end{document}
